% GBEK Non-Equilibrium DMFT: Implementation Status
%
% Reproducible build: run ./fetch_arxiv_figures.sh and ./regenerate_plots.sh
% first (or ensure their outputs already exist alongside this file), then
%   latexmk -pdf report.tex
% Neither the paper figures (arxiv_figures/) nor our own plots (*.png) are
% checked into the repo -- see .gitignore and each script's own header.
%
% \PIVERSION toggles the "bug found and fixed" development-history callouts
% and the bug-count bullet in "Where this leaves us" -- everything else
% (status, capabilities, limitations) is unaffected. Default (0) is the
% full internal version. For a status-only copy (e.g. to share with people
% who aren't interested in development-process detail), build
% report-summary.tex instead: `latexmk -pdf report-summary.tex`.
\documentclass[11pt]{article}
\providecommand{\PIVERSION}{0}
\newif\ifdevnotes
\ifnum\PIVERSION=0 \devnotestrue \else \devnotesfalse \fi

\usepackage[margin=1in]{geometry}
\usepackage{graphicx}
\usepackage{booktabs}
\usepackage{float}
\usepackage[section]{placeins}
\usepackage{amsmath}
\usepackage{amssymb}
\usepackage{siunitx}
\usepackage[colorlinks=true,linkcolor=blue,citecolor=blue,urlcolor=blue]{hyperref}
\usepackage[most]{tcolorbox}
\usepackage{caption}
\usepackage{subcaption}
\usepackage{parskip}
\usepackage{titlesec}

\titleformat{\section}{\normalfont\Large\bfseries}{\thesection}{1em}{}
\titleformat{\subsection}{\normalfont\large\bfseries}{\thesubsection}{1em}{}

\newtcolorbox{callout}{
  colback=blue!4!white, colframe=blue!25!black!40,
  boxrule=0.4pt, arc=2pt, left=8pt, right=8pt, top=6pt, bottom=6pt
}

\newcommand{\Ekin}{\ensuremath{\langle E_\mathrm{kin}\rangle}}
\newcommand{\Eint}{\ensuremath{\langle E_\mathrm{int}\rangle}}
\newcommand{\Etot}{\ensuremath{\langle E_\mathrm{tot}\rangle}}
\newcommand{\Lp}{\ensuremath{\Lambda_+}}

\title{\Large Reproducing GBEK's hybridization-expansion impurity solver}
\author{}
\date{}

\begin{document}
\maketitle
\vspace{-3.5em}

\begin{center}
\begin{minipage}{0.85\textwidth}
\small
\emph{Reference: Gramsch, Balzer, Eckstein \& Kollar, ``Hamiltonian-based
impurity solver for nonequilibrium dynamical mean-field theory,'' Phys.\
Rev.\ B \textbf{88}, 235106 (2013), Sec.\ VI
(\href{https://arxiv.org/abs/1306.6315}{arXiv:1306.6315}).}

\medskip
An independent, from-scratch validation of the paper's low-rank Cholesky /
eigenvector hybridization decomposition and its atomic-limit DMFT
self-consistency scheme --- what we built, how closely it reproduces the
paper's own Figs.\ 3, 4, 7, 8, 9, and 10, and how the C++ production solver
now compares directly against the Python reference for Figs.\ 9--10.
\end{minipage}
\end{center}

\vspace{1.5em}

\section*{What we've implemented}

\begin{itemize}
  \item \textbf{Standalone Python reference.} An independent
    exact-diagonalization implementation of the paper's atomic-limit
    star-geometry SIAM (\texttt{gbek\_reference/}), with no dependency on
    the production C++ solver --- used as ground truth throughout.
  \item \textbf{Both decomposition schemes.} The causal, row-by-row
    \textbf{Cholesky} decomposition of the Weiss field (paper Eqs.\
    56--63) and the non-causal, globally-optimal \textbf{eigenvector}
    decomposition, both as drop-in low-rank factorizations.
  \item \textbf{Full DMFT self-consistency loop.} The fixed-$t_\text{max}$
    iteration scheme (paper Fig.\ 6): hopping-quench ramp, Bethe-lattice
    self-consistency, and iteration to convergence for either
    decomposition mode.
  \item \textbf{Double-occupation observable.} $\langle n_\uparrow(t)
    n_\downarrow(t)\rangle$ tracked directly from the propagated
    many-body state, per-iteration and at convergence, for both
    decomposition modes.
  \item \textbf{Production C++ solver.} \texttt{NeqBathDecomposition.h} /
    \texttt{ImpuritySolverNeqGBEK.h}, cross-checked end-to-end against the
    Python reference on real self-consistency runs, not just unit tests.
    A real two-time Green's-function seeding bug was found and fixed this
    way (see callout below).
  \item \textbf{Regression test suite.} Catch2 coverage for the
    decomposition algorithms (seeding, reconstruction, complex-phase
    targets) plus the impurity solver's Hamiltonian construction.
\end{itemize}

\section{Fig.\ 3 --- Low-rank decomposition quality}
\subsection*{Stepwise reconstruction error, atomic limit ($U=2$)}
The paper's own benchmark for how well a fixed-rank causal (Cholesky) vs.\
non-causal (eigenvector) decomposition reconstructs the exact Weiss field,
via the global relative error $\mathrm{err}[A]$ and its time-resolved
analogue $\mathrm{err}^\mathrm{step}(t)$.

\begin{figure}[h]
  \centering
  \includegraphics[width=0.75\textwidth]{arxiv_figures/paper_fig3.png}
  \caption*{\textbf{Paper} --- input Weiss field, rank-2/3 Cholesky and
    eigenvector reconstructions, and the stepwise error
    $\mathrm{err}^\mathrm{step}(t)$ (bottom left).}
\end{figure}

\begin{figure}[h]
  \centering
  \begin{subfigure}{0.48\textwidth}
    \includegraphics[width=\textwidth]{fig3_errstep.png}
    \caption*{\textbf{Ours} --- $\mathrm{err}^\mathrm{step}(t)$ for rank
      2/3, Cholesky vs.\ eigenvector --- directly comparable to the
      paper's bottom-left panel above.}
  \end{subfigure}
  \hfill
  \begin{subfigure}{0.48\textwidth}
    \includegraphics[width=\textwidth]{gbek_reference_comparison.png}
    \caption*{\textbf{Ours} --- exact reference vs.\ rank-3 Cholesky
      reconstruction: 2D map (paper's own colormap), residual, and
      fixed-$t'$ slices.}
  \end{subfigure}
\end{figure}

\begin{table}[h]
  \centering
  \begin{tabular}{lccc}
    \toprule
    Quantity & Paper (rank 3) & Ours (rank 3) & \\
    \midrule
    err[Cholesky] & 0.17 & 0.168 & match \\
    err[eigenvector] & 0.09 & 0.090 & match \\
    Cholesky / eigenvector ratio & 1.9$\times$ & 1.87$\times$ & match \\
    \bottomrule
  \end{tabular}
\end{table}

\ifdevnotes
\begin{callout}
\textbf{Bug found and fixed (2026-07-13):} the production solver was
seeding $c_{\mathrm{imp},\uparrow}$ ONCE at $t=0$ and propagating the
resulting sector state forward using only that sector's own Hamiltonian
--- a shortcut that is only valid if $c$ commutes with $H$, which the
Hubbard-$U$ term violates. The correct Heisenberg-picture construction
re-seeds $c_{\mathrm{imp},\uparrow}$/$c^\dagger_{\mathrm{imp},\uparrow}$
against the propagated $N$-sector reference state at \emph{every} time
step. The bug produced a systematically suppressed off-diagonal (two-time)
imaginary part, growing with $|n-j|$ --- confirmed independently against a
from-scratch, brute-force dense-matrix-exponential ground truth
(\texttt{cross\_check\_seed\_scheme.py}) before the fix, and reconfirmed
to agree with that same ground truth to machine precision afterward. A
permanent Catch2 regression test now drives the fixed construction
directly and checks it against that ground truth.
\end{callout}

\begin{callout}
\textbf{Bug found and fixed (2026-07-14):} the second-bath (Cholesky)
coupling gave the occupied and empty auxiliary orbitals the \emph{same}
$V(t)$/$V(t)^*$ polarity. GBEK Eq.\ 49/52a requires them to be
\emph{opposite} ($-i\Lambda^<_+$ couples the occupied orbital via
$V(t)V(t')^*$, while $+i\Lambda^>_+ = (-i\Lambda^<_+)^*$ couples the empty
orbital) --- invisible for real $V$, but breaking particle--hole symmetry
as soon as $V$ is complex, which it generally is. Localized by an
independent Python cross-check proving $G^\alpha_\uparrow(t,t) +
G^\beta_\uparrow(t,t) = 1$ exactly, for any bath $V(t)$, when the
$G^\alpha$/$G^\beta$ averaging (Eq.\ 70) is implemented correctly;
cincuenta's own sums were instead 1.01--1.05. All downstream dumps,
figures, and metrics on this page were regenerated after the fix.
\end{callout}
\fi

\section{Fig.\ 4 --- Hybridization amplitude time evolution}
\subsection*{$V_{0p}(t)$, rank $L=3$, Cholesky vs.\ eigenvector}
The paper's illustration of WHY causality matters at the level of the
actual hybridization amplitudes (not just the reconstruction error):
Cholesky's $V^\mathrm{ch}_{0p}(t)$ starts exactly at 0 (since
$\Lambda_+(0,t)=0$), while the eigenvector decomposition's
$V^\mathrm{ev}_{0p}(t)$ does not respect that structure at all.

\begin{figure}[h]
  \centering
  \begin{subfigure}{0.48\textwidth}
    \includegraphics[width=\textwidth]{arxiv_figures/paper_fig4.png}
    \caption*{\textbf{Paper} --- $V^\mathrm{ch}_{0p}(t)$ (top),
      $V^\mathrm{ev}_{0p}(t)$ (bottom), rank $L=3$, with eigenvalue-decay
      inset.}
  \end{subfigure}
  \hfill
  \begin{subfigure}{0.48\textwidth}
    \includegraphics[width=\textwidth]{fig4_hybridization.png}
    \caption*{\textbf{Ours} --- same rank-3 decomposition of our
      atomic-limit target; dashed = imaginary part.}
  \end{subfigure}
\end{figure}

\begin{callout}
\textbf{$V^\mathrm{ch}_{0p}(0)=0$ exactly} in our reproduction (all three
ranks), matching the paper's exact structural claim
($\Lambda_+(0,t)=\Lambda_+(t,0)=0$ forces this). Eigenvalue decay of
$-i\Lambda^<_+$ is rapid in both, qualitatively matching the paper's inset
(roughly 3 orders of magnitude by $p=8$, vs.\ the paper's similar decay by
$p=12$ --- not pixel-matched, but the same qualitative rapid-decay
structure that justifies truncating to a low rank at all).
\end{callout}

\section{Fig.\ 7 --- Causal convergence across DMFT iterations}
\subsection*{Double occupation per iteration ($U=5$, $L_\mathrm{bath}=4$)}
The paper's central structural claim: Cholesky decomposition is
\emph{causal} --- each DMFT iteration only changes the solution beyond a
growing time $t^*$, so results below $t^*$ are locked in early.
Eigenvector decomposition has no such property.

\begin{figure}[h]
  \centering
  \begin{subfigure}{0.48\textwidth}
    \includegraphics[width=\textwidth]{arxiv_figures/paper_fig7.png}
    \caption*{\textbf{Paper} --- seven DMFT iterations, eigenvector (top)
      vs.\ Cholesky (bottom); final iteration solid red.}
  \end{subfigure}
  \hfill
  \begin{subfigure}{0.48\textwidth}
    \includegraphics[width=\textwidth]{fig7_docc_L2.png}
    \caption*{\textbf{Ours} --- same protocol, $L=2$ ($L_\mathrm{bath}=4$,
      matching the paper's own case at left). Cholesky iterations
      converge tightly; eigenvector does not.}
  \end{subfigure}
\end{figure}

\begin{figure}[h]
  \centering
  \includegraphics[width=0.75\textwidth]{fig7_docc_L4.png}
  \caption*{\textbf{Ours, extra} --- $L=4$ ($L_\mathrm{bath}=8$), same
    protocol --- not in the paper's own Fig.\ 7, included here to show the
    causal-convergence story holds at a second rank too.}
\end{figure}

\begin{callout}
Our Cholesky iterations reach a \textbf{non-decreasing, growing agreement
window} that hits $t^* = t_\mathrm{max}$ exactly by the seventh iteration
at both ranks tested --- the paper's convergence criterion, met cleanly.
Eigenvector iterations show no such window at either rank, matching the
paper's claim that it lacks the causal structure entirely.
\end{callout}

\section{Fig.\ 8 --- Converged results vs.\ bath size}
\subsection*{Double occupation at convergence, across ($L_\mathrm{bath}$,
$t_\mathrm{max}$)}
The paper's payoff for causality: a converged Cholesky run's answer for
$t < t^*$ doesn't change when $t_\mathrm{max}$ is extended, so it's safe
to trust up to a horizon that grows with $L_\mathrm{bath}$. Eigenvector
decomposition can converge to a visibly wrong answer instead.

\begin{figure}[h]
  \centering
  \begin{subfigure}{0.48\textwidth}
    \includegraphics[width=\textwidth]{arxiv_figures/paper_fig8.png}
    \caption*{\textbf{Paper} --- four ($L_\mathrm{bath}$,
      $t_\mathrm{max}$) combinations, eigenvector (top) vs.\ Cholesky
      (bottom).}
  \end{subfigure}
  \hfill
  \begin{subfigure}{0.48\textwidth}
    \includegraphics[width=\textwidth]{fig8_docc.png}
    \caption*{\textbf{Ours} --- same four combinations, plus
      $L_\mathrm{bath}=8$ at $t_\mathrm{max}=4$ as an extra rank check.}
  \end{subfigure}
\end{figure}

\begin{table}[h]
  \centering
  \begin{tabular}{lccc}
    \toprule
    Comparison & Paper & Ours & \\
    \midrule
    $L_\mathrm{bath}=4$ vs.\ 8 agreement horizon & $t^*\approx 1.2$--$1.3$ & $t^*\approx 1.16$--$1.2$ & match \\
    Same $L_\mathrm{bath}$, different $t_\mathrm{max}$ & unaffected below $t^*$ & unaffected below $t^*$ & match \\
    \bottomrule
  \end{tabular}
\end{table}

\section{Fig.\ 9 --- Test of energy conservation}
\subsection*{\Ekin, \Eint, \Etot\ vs.\ time, $U=2$ and $U=4$}
The paper's numerical sanity check: once the ramp completes at $t_q$,
\Etot\ should stay flat at its atomic-limit value $-U/4$, with the
residual drift shrinking as $L_\mathrm{bath}$ grows $4\to6\to8$.

\begin{figure}[h]
  \centering
  \begin{subfigure}{0.48\textwidth}
    \includegraphics[width=\textwidth]{arxiv_figures/paper_fig9.png}
    \caption*{\textbf{Paper} --- thick = $L_\mathrm{bath}=8$, thinner = 6,
      4; solid = $U=2$, dashed = $U=4$.}
  \end{subfigure}
  \hfill
  \begin{subfigure}{0.48\textwidth}
    \includegraphics[width=\textwidth]{fig9_energy.png}
    \caption*{\textbf{Ours} --- same protocol and layout, $L=2/3/4$
      ($L_\mathrm{bath}=4/6/8$).}
  \end{subfigure}
\end{figure}

\begin{figure}[h]
  \centering
  \includegraphics[width=0.9\textwidth]{cpp_vs_python_fig9.png}
  \caption*{\textbf{Ours, C++ vs.\ Python} --- same protocol, all three
    $L_\mathrm{bath}$, Python reference (left) vs.\ production C++ solver
    (right) --- visually indistinguishable.}
\end{figure}

\begin{table}[h]
  \centering
  \begin{tabular}{cccc}
    \toprule
    $U$ & $L_\mathrm{bath}$ & $\max|\Etot(t)-\Etot(t_q)|$ for $t>t_q$ (Python) & $\max|\Etot^\mathrm{C++}-\Etot^\mathrm{Python}|$ \\
    \midrule
    2 & 4 & 0.291 & 0.0131 \\
    2 & 6 & 0.197 & 0.0117 \\
    2 & 8 & 0.038 & 0.0117 \\
    4 & 4 & 0.138 & 0.0179 \\
    4 & 6 & 0.134 & 0.0163 \\
    4 & 8 & 0.125 & 0.0163 \\
    \bottomrule
  \end{tabular}
\end{table}

\ifdevnotes
\begin{callout}
\textbf{Bug found and fixed (2026-07-15):} the direct $\langle
H_\mathrm{hop}(t)\rangle$ expectation value used for \Ekin\ came out at
exactly \textbf{2$\times$} the physical kinetic energy. Each bath pair
couples to the impurity through two auxiliary sites --- an ``occ''
partner carrying $\Lambda^<$ and an ``empty'' partner carrying
$\Lambda^>$ --- each with the same coupling magnitude $|V_p(t)|$; correct
for reproducing $\Lambda_+$ itself (confirmed elsewhere: $d(t)$ matches
Fig.\ 10 exactly) but it double-counts the single physical hybridization
channel when read directly off $\langle H_\mathrm{hop}\rangle$. Caught
via the exact conservation law $\Etot=-U/4$: the raw value badly violated
it immediately after the ramp (not just at late times) for the
best-converged $L_\mathrm{bath}=8$ case, while halving $E_\mathrm{kin}$
restored the paper's own pattern --- flat to a few percent at
$L_\mathrm{bath}=8$, with the violation shrinking monotonically as
$L_\mathrm{bath}$ grows. Fixed with a documented $0.5\times$ factor in
\texttt{compute\_energy\_observables} (\texttt{gbek\_selfconsistency.py});
all Fig.\ 9 data and the plot above were regenerated after the fix.
\end{callout}

\begin{callout}
\textbf{C++ port and a second bug (2026-07-15/16):} ported the same
$d(t)$/\Ekin/\Eint/\Etot\ observables into the production C++ solver
(\texttt{ImpuritySolverNeqGBEK::dumpDoccAndEnergy}), including the same
$0.5\times$ factor. An overnight verification run then found
$L_\mathrm{bath}=8$ ($L=4$) gave a WRONG $t=0$ state in C++
($d(0)\approx0.24$, $E_\mathrm{kin}(0)\approx0.47$ instead of the exact
0, 0) --- a separate, pre-existing bug: \texttt{lanczosGS()} (used only
when the extended Fock space exceeds 8 sites) seeded the Lanczos
ground-state search with a naive uniform vector, which converged to the
wrong extremal state (reported energy $-2807.96$ vs.\ the true $-4000$
for $L=4$'s fully-polarized product state). Fixed by seeding Lanczos with
the known analytic atomic-limit product state instead --- exact for this
case, since the pre-quench Hamiltonian has no hopping at all. All three
$L_\mathrm{bath}$ values now agree with Python to within 0.012--0.018 in
$E_\mathrm{tot}$, with no exclusions.
\end{callout}
\fi

\section{Fig.\ 10 --- Double occupation vs.\ $U$}
\subsection*{Converged $d(t)$ for $L_\mathrm{bath}=4$, 6, 8 across a $U$-scan}
Converged double occupation for $U = 0, 1, 2, 4, 6, 8$ --- increasing
suppression with $U$, and a shrinking maximum reliable time as $U$ grows,
exactly as the paper describes.

\begin{figure}[h]
  \centering
  \begin{subfigure}{0.48\textwidth}
    \includegraphics[width=\textwidth]{arxiv_figures/paper_fig10.png}
    \caption*{\textbf{Paper} --- $d(t)$ for $U=0,1,2,4,6,8$,
      $L_\mathrm{bath}=4/6/8$.}
  \end{subfigure}
  \hfill
  \begin{subfigure}{0.48\textwidth}
    \includegraphics[width=\textwidth]{fig10_docc.png}
    \caption*{\textbf{Ours} --- same $U$/$L_\mathrm{bath}$ grid, one panel
      per $L_\mathrm{bath}$.}
  \end{subfigure}
\end{figure}

\begin{callout}
\textbf{$U=0$ closed-form anchor:} the paper states $\Etot=\Ekin=0$ and
$d(t)\to0.25$ exactly for $U=0$. Our $U=0$ runs land at $|E| <
6\times10^{-7}$ and $d(t\to t_\mathrm{max})$ matches 0.25 to plotting
precision at all three $L_\mathrm{bath}$\ifdevnotes, both before and
after the Fig.\ 9 fix (this anchor can't catch an overall scale error
since it's identically zero either way --- see the Fig.\ 9 callout for
how the $E_\mathrm{kin}$ bug was actually caught)\fi. No bugs found in
the $d(t)$ pipeline itself; suppression with $U$ and the shrinking
reliable-time-window with $U$ both match the paper's qualitative
description.
\end{callout}

\begin{figure}[h]
  \centering
  \includegraphics[width=0.9\textwidth]{cpp_vs_python_fig10.png}
  \caption*{\textbf{Ours, C++ vs.\ Python} --- $d(t)$, Python (solid)
    overlaid with C++ (dashed) for the $U=0,2,4$ subset the C++ side has
    data for, one panel per $L_\mathrm{bath}$. Pairs are visually on top
    of each other; $\max|d^\mathrm{C++}-d^\mathrm{Python}|$ is
    0.0038--0.0062 across all nine ($U$, $L_\mathrm{bath}$)
    combinations.}
\end{figure}

\section{Current scope \& limitations}
What the current implementation can and cannot do, and where the
practical boundaries actually sit --- not every constraint here is a bug
to fix; some are deliberate scope choices for this stage of the work.

\begin{table}[H]
  \centering
  \small
  \begin{tabular}{p{3.1cm}p{3.4cm}p{7.5cm}}
    \toprule
    Constraint & Current state & Where it lives \\
    \midrule
    Filling &
    Half-filling only &
    General path requires $n_\uparrow+n_\downarrow=n_\mathrm{sites}$;
    atomic-limit bypass requires $n_\uparrow+n_\downarrow=1$ (single
    spin-polarized electron). $\mu_\mathrm{imp}=-U/2$ is hardcoded for
    particle--hole symmetry in both. Spin-imbalanced fillings
    ($n_\uparrow\neq n_\downarrow$) at half filling overall ARE
    supported. \\
    \addlinespace
    Quench type &
    2 mechanisms, not fully composable &
    Interaction quench $U_i\to U_f$ (instantaneous step only, no ramp
    shape) is available in the general ($n_\mathrm{bath}>0$) path.
    Hopping/bandwidth quench $v(t)$ supports step, cosine, or tanh ramps
    (\texttt{QuenchShape}/\texttt{QuenchDuration}). The atomic-limit
    bypass (\texttt{NeqAtomicLimit=1}) additionally requires $U$ fixed
    ($U_i=U_f$) --- hopping quench only, no interaction quench. \\
    \addlinespace
    First (equilibrium) bath &
    Fixed, not itself quenched &
    The first bath's parameters are set once from the equilibrium fit
    and held fixed through the whole real-time run; only the second
    (GBEK Cholesky) bath captures genuinely non-equilibrium/transient
    hybridization. \\
    \addlinespace
    Ground-state solver switch &
    $n_\mathrm{sites,ext}\le8\to$ dense; $>8\to$ Lanczos &
    $n_\mathrm{sites,ext}=n_\mathrm{bath}+1+2L$. Dense diagonalization is
    exact but $O(\mathrm{dim}^3)$ time / $O(\mathrm{dim}^2)$ memory ---
    infeasible well before $L=5$ ($\mathrm{dim}\sim213$k). For the
    atomic limit ($n_\mathrm{bath}=0$) this puts the switch exactly at
    $L=4$. Lanczos ground-state seeding at $L\ge4$ is now verified
    correct (checked against the exact $t=0$ atomic-limit state) through
    $L=4$. \\
    \addlinespace
    Second-bath rank / max $L_\mathrm{bath}$ tested &
    Verified through $L=4$ ($L_\mathrm{bath}=8$) &
    An $L=5$ ($L_\mathrm{bath}=10$) atomic-limit input exists
    (\texttt{inputNeqAtomicLimitGBEKL5.ain}) but is untested this
    session --- $\mathrm{dim}\sim213$k there, Lanczos-only (dense is not
    an option). \\
    \addlinespace
    Impurity orbitals &
    Single orbital only &
    No multi-orbital Hubbard impurity; the whole GBEK scheme as
    implemented here is built around one impurity site. \\
    \addlinespace
    Known open bug (unrelated) &
    Krylov diagonalization crash &
    \texttt{inputNeqGBEKFig3L3.ain} (the near-atomic-limit approximation
    with a small real first bath, NOT the \texttt{NeqAtomicLimit=1}
    bypass) crashes at step 3/100 post the occ/empty conjugation fix ---
    open, not yet investigated, unrelated to the lanczosGS fix. \\
    \bottomrule
  \end{tabular}
\end{table}

\section*{Where this leaves us}

\begin{itemize}
  \item \textbf{Figs.\ 3, 4, 7, 8, 9, 10} --- all reproduced to close
    quantitative/qualitative agreement with the paper. Figs.\ 9--10 now
    have a direct C++-vs-Python comparison too, not just
    Python-reference-only.
  \ifdevnotes
  \item \textbf{Eight bugs found and fixed} across this effort, all
    sign/conjugation/normalization/seeding mistakes in two-time
    bookkeeping, Hamiltonian construction, or ground-state finding ---
    three in the production C++ solver (two Fig.\ 3 callouts, plus the
    lanczosGS seeding bug in the Fig.\ 9 callout) and, in the Python
    reference, a factor-of-2 in the direct kinetic-energy expectation
    value. Every one was caught by an independent cross-check (a
    closed-form limit, a conservation law, a from-scratch brute-force
    reimplementation, or an exact boundary condition), not by
    inspection.
  \else
  \item \textbf{Validated against the paper at every stage} --- every
    reproduced figure was checked against independent cross-checks
    (closed-form limits, conservation laws, brute-force reference
    calculations, exact boundary conditions), not by inspection alone.
  \fi
  \item \textbf{C++ port of Figs.\ 9--10 done.}
    \texttt{ImpuritySolverNeqGBEK::dumpDoccAndEnergy} now computes and
    dumps the same observables the production solver's Green's functions
    get, unconditionally, alongside them. Direct route ($\langle\psi|
    H_\mathrm{hop}(t)|\psi\rangle$), not the Keldysh convolution the
    Python side abandoned --- matches Python to 0.004--0.018 across all
    tested ($U$, $L_\mathrm{bath}$) combinations, all three
    $L_\mathrm{bath}$ values, no exclusions.
  \item \textbf{C++ is not less efficient than Python for matched
    problems} --- an early apples-to-oranges comparison (a much bigger
    Hilbert space on the C++ side) looked alarming, but timed properly
    against Python's own atomic-limit inputs, C++ was 16$\times$ faster
    at $L_\mathrm{bath}=4$, 2.5$\times$ faster at $L_\mathrm{bath}=6$,
    1.3$\times$ faster at $L_\mathrm{bath}=8$. The margin shrinks with
    $L$ but never reverses in this range.
  \item \textbf{C++ Fig.\ 3 parity} --- the production solver's own
    online reconstruction now overlays the Python Cholesky reference
    essentially exactly on full self-consistency runs\ifdevnotes\ (after
    fixing the two-time Green's-function seeding bug, eliminating a
    spurious ``wing'' artifact that used to appear in the reconstructed
    Weiss field)\fi.
  \item \textbf{Next}: investigate the open $L=3$ near-atomic Krylov
    crash (see the limitations table above); consider testing $L=5$ now
    that Lanczos seeding is fixed; arbitrary (non-half) filling remains
    unimplemented.
\end{itemize}

\vfill
{\small\it cincuenta / gbek\_reference --- DMRG++ --- generated from the
current state of the codebase.}

\end{document}
